\documentclass[10pt,a4paper]{amsart}
\usepackage[margin=19mm]{geometry}
\usepackage{amsmath,amssymb,mathtools}
\usepackage[scr=rsfs,cal=euler]{mathalfa}
\usepackage{xfrac}
\usepackage[unicode,hidelinks,bookmarks=false]{hyperref}
\newcommand{\ie}{i.e.\ }
\newcommand{\cf}{cf.\ }
\newcommand{\resp}{respectively\ }
\newcommand{\Subsection}{\subsection}
\providecommand{\nobreakdash}{\nobreak\hbox{-}}
\setlength{\emergencystretch}{2em}
\multlinegap0pt
\usepackage{amssymb}
\usepackage{mathtools}
\usepackage{csquotes}
\usepackage[shortlabels]{enumitem}
\newtheorem{theorem}{Theorem}
\newtheorem{corollary}{Corollary}
\newtheorem{lemma}{Lemma}
\newcommand{\Gcl}{\mathcal{G}}
\newcommand{\Lcl}{\mathcal{L}}
\newcommand{\SemCone}{\mathcal{C}_\mathrm{sem}}
\title{Strong edge-colouring via local flag algebras}
\author{Eoin Davey, Eoin Hurley, R\'emi de Joannis de Verclos, Ross J. Kang, Jan Volec}
\date{Source: arXiv:2607.17421v1 (CC BY 4.0)}
\begin{document}
\maketitle

\noindent\emph{Source and licence.} Strong edge-colouring via local flag algebras. Version arXiv:2607.17421v1 (CC BY 4.0), distributed by arXiv under the Creative Commons Attribution 4.0 International licence, http://creativecommons.org/licenses/by/4.0/, as recorded in the arXiv licence metadata for this exact version and shown on the abstract page https://arxiv.org/abs/2607.17421v1. Full licence notice: \texttt{licenses/arxiv-2607.17421-CC-BY-4.0.txt}.\par\smallskip

\noindent\textit{Editorial scope.} Counted unit: the article's theorem thm:sec-general (source0.tex lines 126--132). The article's complete original proof of it is delivered with it (source0.tex lines 587--636, one \texttt{proof} environment of 50 lines, which the article places away from the statement). No mathematics is written, repaired or re-derived: the statement and the proof are the article's own lines, in the article's own notation, and no retained line is substituted or re-flowed.  Retained uncounted as the source-local context the proof needs: lines 359--366; lines 403--408; lines 421--429; lines 503--507; lines 566--575.  Nothing was omitted from any retained span, so the package carries no omission record.  No retained line contains a citation, so the wrapper carries no bibliography and no bibliography entry.  No retained line contains a figure environment, an image-inclusion command or drawing code, so no image is delivered and the package declares no asset dependency.  Editorial formatting only, in addition: an AMS \texttt{amsart} preamble that restates the article's own declarations and macros used by the retained lines, namely the environments \texttt{theorem}, \texttt{corollary}, \texttt{lemma} on separate counters, together with the standard packages those lines need.  The retained environments print in this document as Theorem 1, Corollary 1, Lemma 1 in order of appearance, whereas the article prints the same items with the numbering of its own full document.  The complete native source of the article as distributed in the arXiv e-print for this exact version is bundled unchanged as \texttt{sources/205571/source0.tex}.
\begin{corollary}[degree-scale form]
\label{cor:hjk-scale}
For every $\varsigma \in (0, 1]$ and $\iota > 0$ there is a
$\Delta_0$ such that the following holds for every $D \geq \Delta_0$.
If $H$ is a graph with $\Delta(H) \leq D$ and
$|E(H[N_H(v)])| \leq (1 - \varsigma)\binom{D}{2}$ for every
$v \in V(H)$, then $\chi(H) \leq (1 - \varepsilon(\varsigma) + \iota)\,D$.
\end{corollary}

\begin{equation}
\label{eq:criticality-lift}
\chi(L(G)^2)
\leq \max\bigl(\chi(L(G)^2[F]),\,
\lceil(2 - \eta)\Delta(G)^2\rceil + 1\bigr).
\end{equation}

\begin{lemma}[regular suffices]
\label{lem:wlog-regular}
Let $c > 0$. If there is a $\Delta_0$ such that every
$\Delta$-regular graph $G'$ with $\Delta(G') \geq \Delta_0$ satisfies
$\chi'_s(G') \leq c\,\Delta(G')^2$, then there is a $\Delta_1$ such
that every graph $G$ with $\Delta(G) \geq \Delta_1$ satisfies
$\chi'_s(G) \leq c\,\Delta(G)^2$. The same holds with the additional
hypothesis ``bipartite'' on both sides.
\end{lemma}

\begin{equation}
\label{eq:EO-bridge}
|E(L(G)^2[N_{L(G)^2[F]}(f)])|
= |E_O(G')| + o(\Delta(G)^4).
\end{equation}

\begin{lemma}[general strong edge-colouring density bound]
\label{lem:sec-general-density}
\[
\tfrac{10644}{1000}\,\emptyset - O
\in \SemCone^\emptyset
\quad\text{in }\Lcl^\emptyset_5,
\]
where $\Gcl$ is the class above. In particular,
$\phi(O) \leq 10.644$ for every $\phi \in \Phi^\emptyset$.
\end{lemma}

\begin{theorem}[general strong chromatic index bound]
\label{thm:sec-general}
For every graph $G$ with $\Delta(G)$ sufficiently large,
\[
\chi'_s(G) \leq 1.73\,\Delta(G)^2.
\]
\end{theorem}

\begin{proof}[Proof of Theorem~\ref{thm:sec-general}]
By Lemma~\ref{lem:wlog-regular} it suffices to bound $\chi'_s(G)$ on
$\Delta$-regular graphs $G$ with $\Delta$ sufficiently large. Fix
$\eta := 0.2703$ and a maximal $F \subseteq V(L(G)^2)$ on which
$L(G)^2[F]$ has minimum degree at least $(2 - \eta)\Delta(G)^2$. For
an edge $f \in F$ let $G'$ be the $(2,2)$-coloured reduction at
$(G, F, f)$; by construction $G' \in \Gcl$.

Let $(G'_k)$ be any $\Delta$-increasing sequence in $\Gcl$ converging
to a limit functional $\phi \in \Phi^\emptyset$ whose objective density
agrees with $\rho(O; G')$ up to $o(1)$. By
Lemma~\ref{lem:sec-general-density}, $\phi(O) \leq 10.644$. Convert
back to an edge count by~\eqref{eq:EO-bridge}: since $\rho(O; G') =
2|E_O(G')|/\binom{\Delta}{2}^2 + o(1)$,
\[
\frac{|E_O(G')|}{\Delta(G)^4}
\leq \frac{10.644}{8} + o(1)
= 1.3305 + o(1),
\]
and~\eqref{eq:EO-bridge} gives
$|E(L(G)^2[N_{L(G)^2[F]}(f)])| \leq 1.3305\,\Delta(G)^4 +
o(\Delta(G)^4)$, uniformly in $f \in F$. Hence $H := L(G)^2[F]$
satisfies the hypotheses of Corollary~\ref{cor:hjk-scale} at the
degree scale $D := 2\Delta(G)^2$: every strong neighbourhood has at
most $2\Delta(G)^2$ edges, so $\Delta(H) \leq D$, and with
$\varsigma := 1 - 10.644/16$ we have $(1 - \varsigma)\binom{D}{2} =
1.3305\,\Delta(G)^4 - O(\Delta(G)^2)$, so the density bound above is
$(1 - \varsigma')$-sparsity at scale $D$ for any fixed
$\varsigma' < \varsigma$ once $\Delta$ is large; we absorb the
difference into $\iota$ below and keep writing $\varsigma$.

Corollary~\ref{cor:hjk-scale} then colours $H$: a short calculation
gives $\varepsilon(\varsigma) = \varsigma/2 - \varsigma^{3/2}/6
\approx 0.135095$, so for any fixed $\iota > 0$,
\[
\chi(L(G)^2[F])
\leq (1 - 0.135095 + \iota)\cdot 2\Delta(G)^2
< (1.72982 + 2\iota)\Delta(G)^2.
\]
Finally the criticality lift~\eqref{eq:criticality-lift} extends the
bound to all of $L(G)^2$: since $2 - \eta = 1.7297$,
\[
\chi(L(G)^2)
\leq \max\bigl((1.72982 + 2\iota)\Delta(G)^2,\,
\lceil 1.7297\,\Delta(G)^2\rceil + 1\bigr),
\]
and for $\iota$ chosen so that $2\iota < 0.00018$ both branches are
strictly below $1.73\Delta(G)^2$ for $\Delta$ large. Hence
$\chi'_s(G) = \chi(L(G)^2) \leq 1.73\,\Delta(G)^2$.
\end{proof}

\end{document}
